\section{Introduction}
\label{sec:introduction}

Dualities are a recurring organizing principle across physics: they connect descriptions that look
different but encode the same underlying content, often exchanging weak and strong control regimes.
The reduced equilibrium state of an open quantum system is generated by the Hamiltonian of mean
force. Starting from a full Hamiltonian
\(H_{\mathrm{tot}} = H_Q + H_X + H_{\mathrm{int}}\), the reduced Gibbs operator is
\(\bar{\rho}_Q(\beta)=\Tr_X e^{-\beta H_{\mathrm{tot}}}\), and the exact mean-force generator is
defined by \(e^{-\beta H_{\mathrm{MF}}}\propto \bar{\rho}_Q\). In finite-coupling regimes this
object exists exactly but is often difficult to use in practice because asymptotic control depends
strongly on which branch of parameter space is being probed. The problem addressed here is therefore
not existence of the mean-force generator, but control-regime transfer: how to convert hard
asymptotics into tractable ones without changing the underlying analytic family
\cite{gelinThermodynamicsSubensembleCanonical2009a,hiltHamiltonianMeanForce2011,talknerColloquiumStatisticalMechanics2020,rivasStrongCouplingThermodynamics2020,burkeStructureHamiltonianMean2024,jarzynskiNonequilibriumWorkTheorem2004,campisiFluctuationTheoremArbitrary2009,trushechkinOpenQuantumSystem2022}.

The key structural move is to factor the reduced operator into a bare thermal part and an influence
part. This makes transparent which ingredients are universal (bare Gibbs geometry) and which carry
bath-dressing information (the influence operator). Once that partition is explicit, two discrete
representative actions become natural rather than ad hoc.

In this context, the most sought-after target is typically the ground-state sector
(\(\beta\to\infty\)), which is often numerically expensive to access directly.

We isolate two discrete structures. First, a strong-weak involution acts on the crossover coordinate
\(y=\log\chi\), exchanging \(\chi\) and \(\chi^{-1}\). Second, an orientation gauge on the
mean-force eigenaxis induces an affine reflection on effective-temperature representatives,
\(\beff' = 2\beta-\beff\). Together they define a duality geometry for exchanging asymptotic
control sectors. Operationally, this maps a hard low-temperature/ground-state computation to a
numerically trivial dual representative regime.

The manuscript is theory-first. Sections 2--5 establish the formal structure independently of
spin-boson details. The spin-boson qubit appears only as Example~1, where closed formulas make the transfer
mechanism explicit and computationally testable.

The main non-claim is pointwise equality under representative inversion. The results concern
asymptotic exchange of control, not parameter-level physical equivalence. A minimal, theory-led
validation protocol is specified later in the manuscript.

Relative to the previous closure manuscript
\cite{mccaulMeanForceHamiltoniansInfluence2026}, this paper shifts emphasis from model-first
derivation to representative-level duality structure, while keeping explicit scope separation from
full spin-boson universality analyses
\cite{leggettDynamicsDissipativeTwostate1987,weissQuantumDissipativeSystems2012,bullaQuantumPhaseTransitionSubohmic2003,winterQuantumPhaseTransitionSubohmic2009}.
